\documentclass[10pt]{article}
 
% Latin Modern: the modern successor to Computer Modern
\usepackage[T1]{fontenc}
\usepackage{lmodern}
\usepackage[utf8]{inputenc}
 
\usepackage[margin=0.5in]{geometry}
\usepackage{parskip}      % space between paragraphs, no indent
\usepackage{enumitem}
\usepackage{setspace}
\usepackage{multicol}
\singlespacing            % single-spaced body
 
\pagestyle{plain}
 

\begin{document}
 
% REMOVE THE PAGE NUMBER
\thispagestyle{empty}

\begin{multicols}{2}
    \noindent
    Daniel Kadyrov\\
    \today\\
    200-Hour Teacher Training 2026-2027\\
    Sonic Yoga NYC
% \vspace{1em}
    \columnbreak
    % align right
    \begin{flushright}
      Date: November 24th, 2026\\
      Day/Time: Tuesday, 12:30 PM\\
      Class: Vinyasa 2 (60 min)\\
      Teacher: Mitch Hipp\\
    \end{flushright}
\end{multicols}


\vspace{1em}
\begin{center}
  \textbf{Yoga Class Observations}
\end{center}

\begin{enumerate}
    \item \textbf{Was there a theme or message for class? If so, how was it carried through?}\\
    This class was leading up to the Thanksgiving holiday so significant attention was made to giving thanks. 

    \item \textbf{Was there Pranayama, Meditation, Mantra or Mudra used?}\\
    Although Mitch never openly cued Ujjayi, many practitioners were actively practcining the breathing technique. Mitch also cued Simha Pranayama (lion's breath). Meditation was providided in the form of Shivasana. The mantra Ohm was initiated at the beginning and end of the class. 

    \item \textbf{Style of sequence: Cycle (waves of asana that repeat and build) or Category (Working on one category of postures at a time) or other?}\\
    The class follows a compartmentalized categorical format featuring cycling within some categories. Mitch splits sun salutations, standing poses, leg balances, twists, and floor poses with an Ashtanga style chataranga flow in between. Each categorical section features a cycle to krama the difficulty up and prepare for the peak pose that climaxes each partition. 
    
    \item \textbf{Was there a peak pose(s)?}\\
    The peak poses included the crow pose, Bakasana, the side crow pose, Parsva Bakasana, and dancer's pose, Natarajasana. 

    \item \textbf{Name the main actions of the peak pose(s):}\\
    Bakasana: Place hands on the mat, bend elbows, lift hips, and balance knees on upper arms. Engage core and gaze forward.\\
    Parsva Bakasana: Place hands on the mat, bend elbows, lift hips, and balance knees on upper arms. Twist torso to one side, engaging core and gazing forward.\\
    Natarajasana: From a standing position, raise the right leg behind the body and grasp the right foot with the right hand. Press the right leg back while pulling the right hand forward, creating a backbend, opening the chest, and enabling the left arm to reach forward as you balance forward on the left leg. 

    \item \textbf{What postures in the sequence prepared the body for the peak pose(s)?}\\
    Mitch prepped for Natarajasana with a variation of the tiger pose, Vyaghrasana, during the warm up sequence and with the tree pose, Vrksasana, during the standing sequence. The crow poses were prepped with a series of plank variations and chataranga flows. The side crow, Parsva Bakasana, was prepped with the earlier crow pose, Bakasana, with chair pose, Utkatasana, and with twists including revolved side angle pose, Parivrtta Parsvakonasana, and twisted chair pose, Parivrtta Utkatasana.

    \item \textbf{Give an example of pose/counterposes in the main sequence and in the finishing sequence.}\\
    In the leg balance sequence, the tree pose, Vrksasana, was followed by a side body stretch on the opposite side.

    \item \textbf{What was the seated/finishing sequence?  Was it related to the standing sequence in any way or a counter energy sequence?}\\
    The final balancing sequence ended with a pigeon pose, Eka Pada Rajakapotasana. Once that was completed, Mitch formally moved into a seated sequence which included seated single leg folds, Janu Sirsasana, which was counterposed with a gate pose, Parighasana. The seated sequence was concluded with a forward fold, Paschimottanasana, which was counterposed with a reverse plank, Purvottanasana. After lowering the class onto their backs, Mitch cued a supine twist, Supta Matsyendrasana, and then Shivasana. The seated sequence was a counter energy sequence to the intense class that preceded it.

    \item \textbf{What did you learn from observing this class?}\\
    As a devout practitioner of Mitch's 75-minute classes, I was able to observe how Mitch condensed his class into a 60-minute format and offer a more accessible version of his class for a smaller group of students, taking into account the afternoon time slot, and matching the energy and skill level of the students in attendance.

\end{enumerate}
\end{document}